\section{Affine Yangian}\label{sec:affine_Yangian}

\subsection{Affine Yangian}

Fix an integer $N \geq 3$ throughout the paper. We use the notation $\{x,y\}=xy+yx$.

\begin{dfn}\label{dfn:Yangian}The affine Yangian $\affY$ is the algebra over $\bbC$ generated by $x_{i,r}^{+}$, $x_{i,r}^{-}$, $h_{i,r}$ ($i \in \mathbb{Z} / N\mathbb{Z}$, $r \in \mathbb{Z}_{\geq 0}$) with parameters $\ve_1, \ve_2 \in \bbC$ subject to the relations:
\begin{gather*}
[h_{i,r}, h_{j,s}] = 0, \qquad \big[x_{i,r}^{+}, x_{j,s}^{-}\big] = \delta_{ij} h_{i, r+s}, \qquad \big[h_{i,0}, x_{j,r}^{\pm}\big] = \pm a_{ij} x_{j,r}^{\pm},\\
\big[h_{i, r+1}, x_{j, s}^{\pm}\big] - \big[h_{i, r}, x_{j, s+1}^{\pm}\big]= \pm a_{ij} \dfrac{\varepsilon_1 + \varepsilon_2}{2} \big\{h_{i, r}, x_{j, s}^{\pm}\big\} - m_{ij} \dfrac{\varepsilon_1 - \varepsilon_2}{2} \big[h_{i, r}, x_{j, s}^{\pm}\big],\\
\big[x_{i, r+1}^{\pm}, x_{j, s}^{\pm}\big] - \big[x_{i, r}^{\pm}, x_{j, s+1}^{\pm}\big]= \pm a_{ij}\dfrac{\varepsilon_1 + \varepsilon_2}{2} \big\{x_{i, r}^{\pm}, x_{j, s}^{\pm}\big\}
- m_{ij} \dfrac{\varepsilon_1 - \varepsilon_2}{2} \big[x_{i, r}^{\pm}, x_{j, s}^{\pm}\big],\\
\sum_{w \in \mathfrak{S}_{1 - a_{ij}}}\big[x_{i,r_{w(1)}}^{\pm}, \big[x_{i,r_{w(2)}}^{\pm}, \dots, \big[x_{i,r_{w(1 - a_{ij})}}^{\pm}, x_{j,s}^{\pm}\big]\dots\big]\big] = 0, \qquad i \neq j,
\end{gather*}
where
\begin{gather*}
a_{ij} =
\begin{cases}
\hphantom{-}2 &\text{if } i=j, \\
-1 &\text{if } i=j \pm 1, \\
\hphantom{-}0 &\text{otherwise,}
\end{cases}\qquad
m_{ij} =
\begin{cases}
\hphantom{-}1 &\text{if } j=i-1, \\
-1 &\text{if } j=i+1, \\
\hphantom{-}0 &\text{otherwise.}
\end{cases}
\end{gather*}
\end{dfn}

\begin{rem}In this paper, we define $\hat\fraksl_N = \fraksl_N \otimes \bbC\big[t,t^{-1}\big] \oplus \bbC c$ without the degree operator. In \cite{MR3861718}, the Yangian $Y_{\hbar}(\frakg)$ of affine type is defined to be an algebra containing the degree opera\-tor~$d$ and one without~$d$ is denoted by~$Y_{\hbar}(\frakg')$. In particular, the algebra defined in Definition~\ref{dfn:Yangian} coincides with $Y_{\hbar,\ve}(\frakg')$ in the notation of \cite[Definition~7.1]{MR3861718}.
\end{rem}

The subalgebra generated by $x_{i,0}^+$, $x_{i,0}^-$, $h_{i,0}$ ($i \in \bbZ / N\bbZ$) is isomorphic to $U\big(\affsl\big)$ (See \cite[Theorem~6.1]{MR2323534} for $N \geq 4$ and \cite[Theorem~6.9]{GRW} in general). We write $x_{i}^{\pm} = x_{i,0}^{\pm}$, $h_{i} = h_{i,0}$ and identify them with the standard Chevalley generators of $\affsl$. Let $\{\alpha_i\}_{i \in \bbZ/N\bbZ}$ be the simple roots of~$\affsl$. The null root $\delta$ is given by $\delta=\sum\limits_{i=0}^{N-1} \alpha_i$. Let $\theta = \sum\limits_{i=1}^{N-1} \alpha_i$ be the highest root of~$\mathfrak{sl}_N$ and~$h_{-\theta}$ the coroot corresponding to $-\theta$. We denote by $\hat{\Delta}$, $\hat{\Delta}_+$, and $\hat{\Delta}_+^{\mathrm{re}}$ the set of roots, positive roots, and positive real roots for $\affsl$, respectively. We need to consider $\affsl \oplus \bbC d$ with the degree operator to deal with the coproduct on~$\affY$. Fix a nondegenerate invariant symmetric bilinear form $(\, ,\, )$ on $\affsl \oplus \bbC d$ such that $(x_i^+,x_i^-)=1$ and denote the induced bilinear form on the dual of the Cartan subalgebra by the same letter.

Let $\hat{W}$ be the Weyl group of $\affsl$ generated by the simple reflections $s_i$ ($i \in \bbZ/N\bbZ$). We denote by~$s_{\theta}$ the reflection corresponding to the highest root $\theta$. For each $\alpha \in \hat{\Delta}$, we assign the translation element $t_{\alpha}$ in $\hat{W}$. For example, we have $s_0 s_{\theta} = t_{\theta}$. The action of $t_{\alpha}$ on the root lattice $Q=\sum\limits_{i=1}^{N-1} \bbZ \alpha_i$ of $\mathfrak{sl}_N$ is given by $t_{\alpha}(\lambda) = \lambda - (\alpha,\lambda)\delta$ for $\lambda \in Q$.

Set $\hbar = \ve_1 + \ve_2$ and $\tilde{h}_{i,1} = h_{i,1} - \frac{\hbar}{2} h_{i}^2$. We can deduce the following identities directly from the defining relations of $\affY$.
\begin{lem}\label{lem:relations}
We have
\begin{gather*}
\big[h_{i,1}, x_{i}^{\pm}\big] = \pm2 x_{i,1}^{\pm} \pm \hbar \big\{h_i,x_i^{\pm}\big\}, \qquad \big[\tilde{h}_{i,1}, x_{i}^{\pm}\big] = \pm2 x_{i,1}^{\pm}, \qquad \big[x_{i,1}^{\pm},x_i^{\pm}\big] = \pm \hbar \big(x_i^{\pm}\big)^2,\\
\big[h_{i,1}, x_{j}^{\pm}\big] = \mp \left( x_{j,1}^{\pm} + \dfrac{\hbar}{2} \big\{h_i,x_j^{\pm}\big\} - m_{ij} \dfrac{\ve_1-\ve_2}{2}x_j^{\pm} \right) \qquad \text{if} \quad i=j \pm 1.
\end{gather*}
\end{lem}

\subsection{Automorphisms}

We introduce algebra (anti-)automorphisms of $\affY$ which will be used later.

Let $\omega$ be the algebra anti-automorphism of $\affY$ defined by $x_{i,r}^{\pm} \mapsto x_{i,r}^{\mp}$, $h_{i,r} \mapsto h_{i,r}$. Here anti-automorphism means that the invertible operator $\omega$ satisfies $\omega(XY)=\omega(Y)\omega(X)$. It is easy to check that it is well-defined. The following algebra automorphism corresponds to a~rotation of the Dynkin diagram.
\begin{prop}[{\cite[Lemma~3.5]{MR2199856}}] The assignment
\begin{gather*}
x_{i,r}^{\pm} \mapsto \sum_{s=0}^r \dbinom{r}{s} \ve_2^{r-s} x_{i-1,s}^{\pm}, \qquad h_{i,r} \mapsto \sum_{s=0}^r \dbinom{r}{s} \ve_2^{r-s} h_{i-1,s}
\end{gather*}
gives an algebra automorphism $\rho$ of $\affY$.
\end{prop}
In particular we have $\rho(x_{i,1}^{\pm}) = x_{i-1,1}^{\pm} + \ve_2 x_{i-1}^{\pm}$.
\begin{rem}The relation between generators $X_{i,r}^{\pm}$, $H_{i,r}$ with parameters $\lambda$, $\beta$ used in~\cite{MR2199856,MR2323534} and ours are as follows:
\begin{gather*}
X_{0,r}^{\pm} = \sum_{s=0}^r \dbinom{r}{s} \left( \dfrac{N}{4}(\ve_1-\ve_2)\right)^{r-s} x_{0,s}^{\pm},\qquad H_{0,r} = \sum_{s=0}^r \dbinom{r}{s} \left( \dfrac{N}{4}(\ve_1-\ve_2)\right)^{r-s} h_{0,s},\\
X_{i,r}^{\pm} = \sum_{s=0}^r \dbinom{r}{s} \left( \dfrac{i}{2}(\ve_1-\ve_2)\right)^{r-s} x_{i,s}^{\pm}, \qquad H_{i,r} = \sum_{s=0}^r \dbinom{r}{s} \left( \dfrac{i}{2}(\ve_1-\ve_2)\right)^{r-s} h_{i,s}, \qquad i \neq 0,\\
\lambda = \hbar, \qquad \beta = \frac{1}{2} \hbar - \dfrac{N}{4}(\ve_1-\ve_2).
\end{gather*}
\end{rem}

\subsection{Coproduct}

A formula for the coproduct on the affine Yangian $\affY$ was stated in \cite{MR2323534}. Guay--Nakajima--Wendlandt~\cite{MR3861718} gave a detailed proof of the well-definedness. To recall it, we consider a bigger algebra $Y\big(\affsl \oplus \bbC d\big)$, which is generated by $x_{i,r}^{\pm}$, $h_{i,r}$, $d$ with defining relations given in \cite[equation~(2.8)]{MR3861718}. Moreover we need certain completions $Y\big(\affsl \oplus \bbC d\big) \widehat{\otimes} Y\big(\affsl \oplus \bbC d\big)$ and $\affY \widehat{\otimes} \affY$ of the tensor products since the coproduct involves infinite sums. See \cite[Section~5]{MR3861718} for the precise definition of the completion.

We define the half Casimir operator $\Omega_+$ as follows. Let $\{u_k\}$ be a $\bbC$-basis of the Cartan subalgebra of $\affsl \oplus \bbC d$ and $\big\{u^k\big\}$ its dual basis with respect to the nondegenerate bilinear form~$(\, ,\, )$. Let $\big\{x_{\alpha}^{(k)}\big\}_{\alpha \in \hat{\Delta}, 1 \leq k \leq \mult \alpha}$ be a root vector basis satisfying $\big(x_{\alpha}^{(k)},x_{-\alpha}^{(l)}\big) = \delta_{k,l}$. Here $\mult \alpha$ denotes the dimension of the root space corresponding to~$\alpha$. We take simple root vectors as $x_{\pm \alpha_i}^{(1)} = x_i^{\pm}$. Put
\begin{gather*}
\Omega_+ = \sum_{k} u^k \otimes u_k + \sum_{\substack{\alpha \in \hat{\Delta}_+ \\ 1 \leq k \leq \mult \alpha}} x_{-\alpha}^{(k)} \otimes x_{\alpha}^{(k)}
\end{gather*}
as an element of $Y\big(\affsl \oplus \bbC d\big) \widehat{\otimes} Y\big(\affsl \oplus \bbC d\big)$.

Define a $\bbC$-linear operator $\square$ on $\affY$ by $\square(X) = X \otimes 1 + 1 \otimes X$. Note that $\square$ is not an algebra homomorphism, but satisfies $\square([X,Y]) = [\square(X), \square(Y)]$.
\begin{thm}[{\cite[Definition~4.6, Theorem~4.9, Proposition~5.18, Section~7]{MR3861718}}]\label{thm:definition_of_coproduct} There exists an algebra homomorphism $\Delta \colon \affY \to \affY \widehat{\otimes} \affY$ uniquely determined by
\begin{gather*}
\Delta(X) = \square(X) \qquad \text{for} \quad X=x_{i}^{\pm}, h_i,\\
\Delta\big(x_{i,1}^+\big) = \square(x_{i,1}^+) - \hbar \big[1 \otimes x_i^+, \Omega_+\big],\\
\Delta\big(x_{i,1}^-\big) = \square(x_{i,1}^-) + \hbar \big[x_i^- \otimes 1, \Omega_+\big],\\
\Delta\big(\tilde{h}_{i,1}\big)= \square(\tilde{h}_{i,1}) + \hbar [h_i \otimes 1, \Omega_+].
\end{gather*}
\end{thm}
\begin{rem}Explicitly we have
\begin{gather*}
[1 \otimes x_i^+, \Omega_+] = -h_i \otimes x_i^+ + \sum_{\substack{\alpha \in \hat{\Delta}_+\\ 1 \leq k \leq \mult \alpha}} x_{-\alpha}^{(k)} \otimes \big[x_i^+, x_{\alpha}^{(k)}\big],
\\
[x_i^- \otimes 1, \Omega_+] = x_i^- \otimes h_i + \sum_{\substack{\alpha \in \hat{\Delta}_+\\ 1 \leq k \leq \mult \alpha}} \big[x_i^-, x_{-\alpha}^{(k)}\big] \otimes x_{\alpha}^{(k)},
\\
[h_i \otimes 1, \Omega_+] = - \sum_{\alpha \in \hat{\Delta}_+^{\mathrm{re}}} (\alpha_i,\alpha)\, x_{-\alpha}^{(1)} \otimes x_{\alpha}^{(1)}.
\end{gather*}
Hence the target of $\Delta$ is $\affY \widehat{\otimes} \affY$ without~$d$.
\end{rem}

\section{Braid group action}\label{section3}

We define automorphisms $T_i$ $(i \in \bbZ/N\bbZ)$ of the affine Yangian and study their properties.

\subsection{Definition}

Since the adjoint actions of $x_i^{\pm}$ on $\affY$ are locally nilpotent, the operators $\exp\ad x_i^{\pm}$ are well defined by
\begin{gather*}
\exp\ad x_i^{\pm} = \sum_{n=0}^{\infty} \dfrac{1}{n!} \big(\ad x_i^{\pm}\big)^n.
\end{gather*}
These are algebra automorphisms of $\affY$ as $\ad x_i^{\pm}$ are derivations. We define an algebra automorphism $T_i$ of $\affY$ for each $i \in \bbZ/N\bbZ$ by
\begin{gather*}
T_i = \exp\ad x_i^{+} \exp\ad (-x_i^{-}) \exp\ad x_i^{+}.
\end{gather*}
This operator appears in \cite{MR3861718} and is used to construct real root vectors of Yangians.

\subsection{Braid relations}

A proof of the following proposition is exactly the same as one for the fact that $\{T_i\}$ satisfy the braid relations as automorphisms of $U\big(\affsl\big)$. We give a proof for the sake of completeness.

\begin{prop}\label{prop:braid}The operators $\{T_i\}$ satisfy the braid relations. That is, we have
\begin{gather*}
T_i T_j = T_j T_i \quad \text{if $a_{ij}=0$},\qquad T_i T_j T_i = T_j T_i T_j \quad \text{if $a_{ij}=-1$}.
\end{gather*}
\end{prop}
For each $w \in \hat{W}$ with a reduced expression $w=s_{i_1} \cdots s_{i_l}$, we can define an algebra automorphism $T_w$ of $\affY$ by $T_w = T_{i_1} \cdots T_{i_l}$ thanks to the braid relations.

Let us start the proof with some preparations. A proof of the following formulas is straightforward.
\begin{lem}\label{lem:expad1} Assume $a_{ij}=-1$. Then
\begin{enumerate}\itemsep=0pt
\item $\exp\ad x_{i}^{+}$ sends:
\begin{gather*}
x_i^+ \mapsto x_i^+, \qquad x_j^+ \mapsto x_j^+ + \big[x_i^+,x_j^+\big],\qquad x_i^- \mapsto x_i^-+h_i-x_i^+, \qquad x_j^- \mapsto x_j^-,\\
h_i \mapsto h_i-2x_i^+,\qquad h_j \mapsto h_j+x_i^+,\qquad \big[x_i^+,x_j^+\big] \mapsto \big[x_i^+,x_j^+\big],\\
 \big[x_i^-,x_j^-\big] \mapsto \big[x_i^-,x_j^-\big] + x_j^-;
\end{gather*}
\item $\exp\ad (-x_{i}^{-})$ sends:
\begin{gather*}
x_i^+ \mapsto x_i^++h_i-x_i^-, \qquad x_j^+ \mapsto x_j^+,\qquad x_i^- \mapsto x_i^-, \qquad x_j^- \mapsto x_j^--[x_i^-,x_j^-],\\
h_i \mapsto h_i-2x_i^-,\qquad h_j \mapsto h_j+x_i^-,\qquad \big[x_i^+,x_j^+\big] \mapsto \big[x_i^+,x_j^+\big]-x_j^+,\\
 \big[x_i^-,x_j^-\big] \mapsto \big[x_i^-,x_j^-\big].
\end{gather*}
\end{enumerate}
\end{lem}

The following two propositions are well known. Proposition~\ref{prop:formula1} follows from Lemma~\ref{lem:expad1}. Then Proposition~\ref{prop:iji_pre} follows from Proposition~\ref{prop:formula1}.
\begin{prop}\label{prop:formula1} We have
\begin{gather*}
T_i(x_j^+) = \begin{cases}
-x_i^- & \text{if $i=j$},\\
\big[x_i^+,x_j^+\big] & \text{if $a_{ij}=-1$},\\
x_j^+ & \text{if $a_{ij}=0$},
\end{cases}\qquad
T_i\big(x_j^-\big) = \begin{cases}
-x_i^+ & \text{if $i=j$},\\
-\big[x_i^-,x_j^-\big] & \text{if $a_{ij}=-1$},\\
x_j^- & \text{if $a_{ij}=0$},
\end{cases}
\\
T_i(h_j) = \begin{cases}
-h_i & \text{if $i=j$},\\
h_j + h_i & \text{if $a_{ij}=-1$},\\
h_j & \text{if $a_{ij}=0$}.
\end{cases}
\end{gather*}
\end{prop}

\begin{prop}\label{prop:iji_pre}Assume $a_{ij}=-1$. Then we have
\begin{gather*}
T_i T_j \big(x_i^+\big) = x_j^+,\qquad T_i T_j \big(x_i^-\big) = x_j^-,\qquad T_i T_j (h_i) = h_j.
\end{gather*}
\end{prop}

\begin{proof}[Proof of Proposition~\ref{prop:braid}] We use that $\varphi \circ \exp\ad X \circ \varphi^{-1} = \exp \ad \varphi(X)$ holds for any algebra automorphism~$\varphi$. If $a_{ij}=0$, then we have
\begin{gather*}
T_i T_j T_i^{-1} = \exp\ad T_i\big(x_j^+\big) \exp\ad T_i\big({-}x_j^-\big) \exp\ad T_i\big(x_j^+\big) = T_j
\end{gather*}
by the formulas in Proposition~\ref{prop:formula1}. If $a_{ij}=-1$, then we have
\begin{gather*}
T_i T_j T_i T_j^{-1} T_i^{-1} = \exp\ad T_i T_j \big(x_i^+\big) \exp\ad T_i T_j \big({-}x_i^-\big) \exp\ad T_i T_j \big(x_i^+\big) = T_j
\end{gather*}
by the formulas in Proposition~\ref{prop:iji_pre}. The proof is complete.
\end{proof}

We give an alternative definition of $T_i$.
\begin{prop}The operator $T_i$ coincides with $\exp\ad \big({-}x_i^{-}\big) \exp\ad x_i^{+} \exp\ad \big({-}x_i^{-}\big)$.
\end{prop}
\begin{proof}Put $T_i' = \exp\ad \big({-}x_i^{-}\big) \exp\ad x_i^{+} \exp\ad \big({-}x_i^{-}\big)$ for a while. Then we have
\begin{gather*}
T_i T_i' T_i^{-1} = \exp\ad T_i\big({-}x_i^{-}\big) \exp\ad T_i\big(x_i^+\big) \exp\ad T_i\big({-}x_i^{-}\big) = T_i
\end{gather*}
by the formulas in Proposition~\ref{prop:formula1}. Hence the assertion is proved.
\end{proof}

Note that $T_i^{-1}$ is given by
\begin{gather*}
\exp\ad x_i^- \exp\ad \big({-}x_i^+\big) \exp\ad x_i^- = \exp\ad \big({-}x_i^+\big) \exp\ad x_i^- \exp\ad \big({-}x_i^+\big).
\end{gather*}

\begin{lem}\label{lem:rho_omega_and_T} We have $(i)$ $\omega \circ T_i = T_{i} \circ \omega$ and $(ii)$ $\rho \circ T_i = T_{i-1} \circ \rho$.
\end{lem}
\begin{proof}For the assertion (i), we use $\omega (\exp\ad X (Y)) = \exp\ad(-\omega(X))(\omega(Y))$. Then we have $\omega(T_i(X))=T_i'(\omega(X))$. Since we have proved $T_i'=T_i$, the assertion holds.

The assertion (ii) is obvious.
\end{proof}

Let $M$ be a $\affY$-module and assume that $x_i^{\pm}$ acts on $M$ locally nilpotently. Then an automorphism $T_i^M$ of $M$ is defined similarly. The following property is immediate.
\begin{prop}Let $M$ be a $\affY$-module and assume that $x_i^{\pm}$ acts on $M$ locally nilpotently. Then for any $X \in \affY$ and $m \in M$,
\begin{gather*}
T_i^M(Xm)=T_i(X)T_i^M(m)
\end{gather*}
holds.
\end{prop}

In fact, this is a general property for any associative algebra such that it contains $U\big(\affsl\big)$ as a subalgebra and the operators $\exp\ad x_i^{\pm}$ are well-defined.

\subsection{Action on generators}

We compute the action of $T_i$ on $x_{j,1}^{\pm}$, $\tilde{h}_{j,1}$. We will use the following formulas.
\begin{lem}\label{lem:expad2}Assume $a_{ij}=-1$. Then
\begin{enumerate}\itemsep=0pt
\item $\exp\ad x_{i}^{+}$ sends:
\begin{gather*}
x_{i,1}^+ \mapsto x_{i,1}^+-\hbar \big(x_i^+\big)^2,\\
x_{i,1}^- \mapsto x_{i,1}^- + h_{i,1}-x_{i,1}^+ -\tfrac{\hbar}{2}\big\{ h_i, x_i^+\big\} + \hbar \big(x_i^+\big)^2,\\
x_{j,1}^- \mapsto x_{j,1}^-,\\
\big[x_i^-,x_{j,1}^-\big] \mapsto \big[x_i^-,x_{j,1}^-\big]+x_{j,1}^-,\\
\big\{h_i,x_i^+\big\} \mapsto \big\{h_i,x_i^+\big\} - 4\big(x_i^+\big)^2,\\
\big(x_i^+\big)^2 \mapsto (x_i^+)^2;
\end{gather*}
\item $\exp\ad (-x_{i}^{-})$ sends:
\begin{gather*}
x_{i,1}^+ \mapsto x_{i,1}^+ + h_{i,1}-x_{i,1}^- -\tfrac{\hbar}{2}\big\{ h_i, x_i^-\big\} + \hbar \big(x_i^-\big)^2,\\
x_{i,1}^- \mapsto x_{i,1}^- - \hbar\big(x_i^-\big)^2,\\
x_{j,1}^- \mapsto x_{j,1}^- - \big[x_i^-,x_{j,1}^-\big],\\
h_{i,1} \mapsto h_{i,1} - 2x_{i,1}^- - \hbar\big\{ h_i, x_i^- \big\} + 3\hbar \big(x_i^-\big)^2,\\
\big\{h_i,x_i^+\big\} \mapsto \big\{h_i,x_i^+\big\} - 3\big\{h_i,x_i^-\big\} - 2\big\{x_i^+,x_i^-\big\} + 2h_i^2 + 4\big(x_i^-\big)^2,\\
\big(x_i^+\big)^2 \mapsto \big(x_i^+\big)^2 + h_i^2 + \big(x_i^-\big)^2+ \big\{h_i,x_i^+\big\} - \big\{h_i,x_i^-\big\} - \big\{x_i^+,x_i^-\big\}.
\end{gather*}
\end{enumerate}
\end{lem}
\begin{proof}By a direct computation using the defining relations of $\affY$ with Lemma~\ref{lem:relations}. For example, by
\begin{gather*}
\ad x_i^{+} \big(x_{i,1}^-\big) = h_{i,1}, \\
\big(\ad x_i^{+}\big)^2 \big(x_{i,1}^-\big) = \big[x_i^+,h_{i,1}\big] = -2x_{i,1}^+ - \hbar \big\{h_i,x_i^+\big\}, \\
\big(\ad x_i^{+}\big)^3 \big(x_{i,1}^-\big) = -2 \big[x_i^+, x_{i,1}^+\big] - \hbar \big\{\big[x_i^+,h_i\big],x_i^+\big\} = 2\hbar\big(x_i^+\big)^2 + 4\hbar \big(x_i^+\big)^2 = 6\hbar\big(x_i^+\big)^2, \\
\big(\ad x_i^{+}\big)^4 \big(x_{i,1}^-\big)=0,
\end{gather*}
we have
\begin{gather*}
\exp\ad x_i^+ \big(x_{i,1}^-\big) = x_{i,1}^- + h_{i,1} - x_{i,1}^+ - \tfrac{\hbar}{2} \big\{h_i,x_i^+\big\} + \hbar\big(x_i^+\big)^2.\tag*{\qed}
\end{gather*}\renewcommand{\qed}{}
\end{proof}


\begin{prop}\label{prop:formula2}We have
\begin{gather*}
T_i\big(x_{j,1}^+\big) = \begin{cases}
-x_{i,1}^- + \tfrac{\hbar}{2}\big\{ h_i, x_i^- \big\} & \text{if $i=j$},\\
\big[x_i^+,x_{j,1}^+\big] & \text{if $a_{ij}=-1$},\\
x_{j,1}^+ & \text{if $a_{ij}=0$},
\end{cases} \\
T_i\big(x_{j,1}^-\big) = \begin{cases}
-x_{i,1}^+ + \tfrac{\hbar}{2}\big\{ h_i, x_i^+ \big\} & \text{if $i=j$},\\
-\big[x_{i}^-,x_{j,1}^-\big] & \text{if $a_{ij}=-1$},\\
x_{j,1}^- & \text{if $a_{ij}=0$},
\end{cases}
\\
T_i\big(\tilde{h}_{j,1}\big) = \begin{cases}
-\tilde{h}_{i,1} - \hbar \big\{ x_i^+,x_i^- \big\} & \text{if $i=j$},\\
\tilde{h}_{j,1} + \tilde{h}_{i,1} + \tfrac{\hbar}{2} \big\{ x_i^+,x_i^- \big\} + m_{ij}\tfrac{\ve_1 - \ve_2}{2} h_i & \text{if $a_{ij}=-1$},\\
\tilde{h}_{j,1} & \text{if $a_{ij}=0$}.
\end{cases}
\end{gather*}
\end{prop}

\begin{proof}The formulas for $a_{ij}=0$ trivially hold. Let us consider the other cases.

We compute $T_i\big(x_{j,1}^-\big)$. First assume $i=j$. We have
\begin{gather*}
\exp\ad \big({-}x_i^-\big)\exp\ad x_i^+\big(x_{i,1}^-\big) = \exp\ad \big({-}x_i^-\big)\big(x_{i,1}^- + h_{i,1}-x_{i,1}^+ -\tfrac{\hbar}{2}\big\{ h_i, x_i^+\big\} + \hbar \big(x_i^+\big)^2\big)\\
\qquad {}= \big(x_{i,1}^- - \hbar(x_i^-)^2\big)+\big(h_{i,1}-2x_{i,1}^- - \hbar\big\{ h_i, x_i^- \big\} + 3\hbar \big(x_i^-\big)^2\big)\\
\qquad\quad{}-\big(x_{i,1}^+ + h_{i,1}-x_{i,1}^--\tfrac{\hbar}{2}\big\{ h_i, x_i^-\big\} + \hbar \big(x_i^-\big)^2\big)\\
\qquad\quad{}-\tfrac{\hbar}{2}\big(\big\{h_i,x_i^+\big\} - 3\big\{h_i,x_i^-\big\} - 2\big\{x_i^+,x_i^-\big\} + 2h_i^2 + 4\big(x_i^-\big)^2\big)\\
\qquad\quad{} +\hbar\big(\big(x_i^+\big)^2 + h_i^2 + \big(x_i^-\big)^2+ \big\{h_i,x_i^+\big\} - \big\{h_i,x_i^-\big\} - \big\{x_i^+,x_i^-\big\}\big)\\
\qquad{}=-x_{i,1}^+ + \tfrac{\hbar}{2}\big\{h_i,x_i^+\big\} +\hbar \big(x_i^+\big)^2.
\end{gather*}
Then we have
\begin{gather*}
\exp\ad x_i^+\big({-}x_{i,1}^+ + \tfrac{\hbar}{2}\big\{h_i,x_i^+\big\} +\hbar \big(x_i^+\big)^2\big) \\
\qquad {}= -\big( x_{i,1}^+-\hbar \big(x_i^+\big)^2 \big) + \tfrac{\hbar}{2}\big( \big\{h_i,x_i^+\big\} - 4\big(x_i^+\big)^2 \big) + \hbar \big(x_i^+\big)^2
 =-x_{i,1}^++\tfrac{\hbar}{2}\big\{h_i,x_i^+\big\}.
\end{gather*}
Next assume $a_{ij}=-1$. Then we have
\begin{gather*}
T_i\big(x_{j,1}^-\big)=\exp\ad x_i^+\exp\ad \big({-}x_i^-\big)\big(x_{j,1}^-\big) = \exp\ad x_i^+\big(x_{j,1}^- - \big[x_i^-,x_{j,1}^-\big]\big)\\
\hphantom{T_i\big(x_{j,1}^-\big)}{} = x_{j,1}^- - \big(\big[x_i^-,x_{j,1}^-\big]+x_{j,1}^-)=-\big[x_i^-,x_{j,1}^-\big]. %\label{eq:formula2-1}
\end{gather*}
The formulas for $T_i\big(x_{j,1}^+\big)$ are obtained by applying $\omega$ to these results.

We compute $T_{i}(h_{j,1})$. First assume $i=j$. Then we have
\begin{gather*}
T_i(h_{i,1}) = T_i\big(\big[x_{i}^+,x_{i,1}^-\big]\big) = \big[{-}x_{i}^-, -x_{i,1}^+ + \tfrac{\hbar}{2}\big\{ h_i, x_i^+ \big\}\big] = -h_{i,1} - \hbar \big\{x_i^+,x_i^-\big\} + \hbar h_i^2.
\end{gather*}
Next assume $a_{ij}=-1$. Then we have
\begin{gather*}
T_i(h_{j,1}) = T_i\big(\big[x_{j}^+,x_{j,1}^-\big]\big) = \big[\big[x_i^+,x_{j}^+\big], -\big[x_i^-,x_{j,1}^-\big]\big] \\
\hphantom{T_i(h_{j,1})}{} = - \big[\big[\big[x_i^+,x_j^+\big],x_i^-\big],x_{j,1}^-\big]-\big[x_i^-,\big[\big[x_i^+,x_j^+\big],x_{j,1}^-\big]\big].
\end{gather*}
Since we have
\begin{gather*}
\big[\big[\big[x_i^+,x_j^+\big],x_i^-\big],x_{j,1}^-\big]=\big[\big[h_i,x_j^+\big],x_{j,1}^-\big]=\big[{-}x_j^+,x_{j,1}^-\big]=-h_{j,1}
\end{gather*}
and
\begin{gather*}
\big[x_i^-,\big[[x_i^+,x_j^+\big],x_{j,1}^-\big]\big] =\big[x_i^-,\big[x_i^+,h_{j,1}\big]\big] = \big[x_i^-,x_{i,1}^++\tfrac{\hbar}{2}\big\{h_j,x_i^+\big\}-m_{ji}\tfrac{\ve_1 - \ve_2}{2} x_i^+\big] \\
\hphantom{\big[x_i^-,\big[[x_i^+,x_j^+\big],x_{j,1}^-\big]\big]}{} = -h_{i,1} + \tfrac{\hbar}{2}\big( \big\{\big[x_i^-,h_j\big],x_i^+\big\} + \big\{ h_j, \big[x_i^-, x_i^+\big] \big\}\big) + m_{ji}\tfrac{\ve_1 - \ve_2}{2} h_i \\
\hphantom{\big[x_i^-,\big[[x_i^+,x_j^+\big],x_{j,1}^-\big]\big]}{} = -h_{i,1} - \tfrac{\hbar}{2} \big\{x_i^-,x_i^+\big\} - \hbar h_i h_j + \tfrac{\ve_1 - \ve_2}{2}m_{ji} h_i,
\end{gather*}
we conclude
\begin{gather*}
T_i(h_{j,1})=h_{j,1}+h_{i,1}+\tfrac{\hbar}{2} \big\{x_i^+,x_i^-\big\} + \hbar h_i h_j + \tfrac{\ve_1 - \ve_2}{2}m_{ij} h_i.
\end{gather*}
We can easily obtain the formulas for $T_{i}(\tilde{h}_{j,1})$ from these identities.
\end{proof}

\begin{prop}\label{prop:iji}Assume $a_{ij}=-1$. Then we have
\begin{gather*}
T_i T_j \big(x_{i,1}^+\big) = x_{j,1}^+ - m_{ij}\tfrac{\ve_1 - \ve_2}{2}x_j^+ - \tfrac{\hbar}{2}\big\{T_i\big(x_j^+\big),x_i^-\big\}, \\
T_i T_j \big(x_{i,1}^-\big) = x_{j,1}^- - m_{ij}\tfrac{\ve_1 - \ve_2}{2}x_j^- - \tfrac{\hbar}{2}\big\{x_i^+,T_i\big(x_j^-\big)\big\}, \\
T_i T_j (h_{i,1}) = h_{j,1} - m_{ij}\tfrac{\ve_1 - \ve_2}{2}h_j - \tfrac{\hbar}{2}\{x_i^+,x_i^-\} + \tfrac{\hbar}{2}\big\{T_i(x_j^+),T_i\big(x_j^-\big)\big\}.
\end{gather*}
\end{prop}
\begin{proof}We show the first identity. Since we have
\begin{gather*}
T_iT_j\big(x_{i,1}^+\big) = T_i\big(\big[x_{j}^+,x_{i,1}^+\big]\big) = \big[\big[x_i^+,x_j^+\big],-x_{i,1}^-+\tfrac{\hbar}{2}\big\{h_i,x_i^{-}\big\}\big],
\end{gather*}
it follows from
\begin{gather*}
\big[\big[x_i^+,x_j^+\big],x_{i,1}^-\big] = \big[h_{i,1},x_j^+\big] = -\big( x_{j,1} + \tfrac{\hbar}{2}\big\{h_i,x_j^{+}\big\} - m_{ij} \tfrac{\ve_1 - \ve_2}{2}x_j^+ \big)
\end{gather*}
and
\begin{gather*}
\big[\big[x_i^+,x_j^+\big],\big\{h_i,x_i^{-}\big\}\big]= \big\{\big[\big[x_i^+,x_j^+\big],h_i\big],x_i^-\big\} + \big\{h_i,\big[\big[x_i^+,x_j^+\big],x_i^-\big]\big\} \\
 \hphantom{\big[\big[x_i^+,x_j^+\big],\big\{h_i,x_i^{-}\big\}\big]}{} = -\big\{\big[x_i^+,x_j^+\big],x_i^-\big\} - \big\{h_i,x_j^{+}\big\} = -\big\{T_i(x_j^+),x_i^-\big\} - \big\{h_i,x_j^{+}\big\}.
\end{gather*}
The second identity is obtained by applying $\omega$ to the first one. The third identity follows from
\begin{gather*}
T_iT_j(h_{i,1})= T_iT_j\big(\big[x_{i,1}^+,x_{i}^-\big]\big) = \big[x_{j,1}^+-m_{ij}\tfrac{\ve_1-\ve_2}{2}x_j^+-\tfrac{\hbar}{2}\big\{T_i\big(x_j^+\big),x_i^-\big\},x_j^-\big]
\end{gather*}
and
\begin{gather*}
\big[\big\{T_i\big(x_j^+\big),x_i^-\big\},x_j^-\big] = \big\{\big[T_i(x_j^+),x_j^-\big],x_i^-\big\} + \big\{T_i\big(x_j^+\big),\big[x_i^-,x_j^-\big]\big\}\\
\hphantom{\big[\big\{T_i\big(x_j^+\big),x_i^-\big\},x_j^-\big]}{} = \big\{x_i^+,x_i^-\big\} - \big\{T_i\big(x_j^+\big),T_i\big(x_j^-\big)\big\}.\tag*{\qed}
\end{gather*}\renewcommand{\qed}{}
\end{proof}

We give formulas for $T_i^2$. Proofs are straightforward.
\begin{prop} We have $T_i^2(X)=X$ for $X=x_{j}^{\pm}, x_{j,1}^{\pm}$ $(a_{ij} \neq -1)$, $h_k, h_{k,1}$ $(k \in \bbZ/N\bbZ)$. We have $T_i^2(X)=-X$ for $X=x_{j}^{\pm}, x_{j,1}^{\pm}$ $(a_{ij} = -1)$. In particular, $T_i^4 = 1$.
\end{prop}

